\section{Limitations}\label{sec:limitations}

While BEACON's design and evaluation framework are grounded in existing research and best practices as much as possible, there are several limitations needing consideration.

First, BEACON is locked into a simulation-only environment.
No physical unmanned aerial vehicle (UAV) or unmanned surface vessel (USV) testing was conducted, and there is no real-time data link between physical drones and the simulation.
This means that BEACON is best described as a Digital Model/Shadow, rather than a full Digital Twin.
Due to this, environmental effects, drone kinematics, communications degradations, and other related factors are simplified, which limits the scope of the conclusions that can be drawn from BEACON\@.
This trade-off was made to focus on human-computer interaction (HCI) and interface design, which is the core contribution of this project, leaving physical testing for future work.

Similarly, human subject testing was also left for future work, due to ethical and logistical concerns with the timeline of this project.
Workload interpretations are based off of theory and informal self-use, rather than formally tested, controlled experiments with human subjects.
This also means that the conclusions drawn from BEACON's evaluation framework are limited, and are not generalizable to real-world operators without further testing.

The simulation also focuses on just searching rather than the full spectrum of MSAR operations, which include tasks like target extraction, first aid, and more.
It also assumes perfect anomaly detection, and does not include coastal features, islands, or anomalies drifting with the current.
These concerns were left for future work to focus on the core mechanics of drone supervision and interface design.
